\section*{Chapter \ref{ch:b3:quantitative-geomorphology} --- Quantitative Geomorphology}

\begin{solution}{exo:b3:quantitative-geomorphology:1}
$a/\sqrt{\pi\kappa t} = 0.30 - 0.04 = 0.26$, so $\kappa t = (2/0.26)^2/\pi =
\qty{18.8}{m^2}$, and $t = 18.8/(5\times10^{-3}) = \qty{3.8}{kyr}$, within a
factor of two either way given the range of $\kappa$.
\end{solution}

\begin{solution}{exo:b3:quantitative-geomorphology:2}
$E = \kappa\,|\nabla^2 z| = 4\times10^{-3}\times0.012 =
\qty{4.8e-5}{m/yr}$, \qty{0.048}{mm/yr}. Relief $E L^2/(2\kappa) =
0.012\times80^2/2 = \qty{38.4}{m}$.
\end{solution}

\begin{solution}{exo:b3:quantitative-geomorphology:3}
$E = 2\times10^{-6}\times(25\times10^6)^{0.5}\times0.01 = 2\times10^{-6}
\times5000\times0.01 = \qty{1.0e-4}{m/yr}$, \qty{0.1}{mm/yr}. Since $E$ is
in \unit{m/yr}, $A^m$ in $\unit{m}^{2m}$ and $S$ has no unit, $K$ is in
$\unit{m}^{1-2m}\,\unit{yr^{-1}}$: \unit{yr^{-1}} for $m = 0.5$.
\end{solution}

\begin{solution}{exo:b3:quantitative-geomorphology:4}
(a) Fixed: a hard bed steepens the river where it crops out, and the step
stays with the bed. (b) Anchored at the fault: as long as the fault slips,
the river is steepened where it crosses it. (c) Mobile: the uplift of the
coast lowered the \omterm{def:b3:quantitative-geomorphology:knickpoint}{base level} of every river at once; the \omterm{def:b3:quantitative-geomorphology:knickpoint}{knickpoints}
are migrating upstream and, as the kinematic-wave theorem predicts for $n
= 1$, they lie on one contour.
\end{solution}

\begin{solution}{exo:b3:quantitative-geomorphology:5}
$\theta = m/n = 0.5$. $S = (U/K)^{1/2}A^{-1/2} = (10^{-4}/10^{-9})^{1/2}
\times(10^7)^{-1/2} = 316\times3.16\times10^{-4} = 0.10$. Doubling $U$
multiplies $k_s = (U/K)^{1/n}$ by $\sqrt2 = 1.41$; for $n = 1$ by 2: the
larger $n$, the less steepness responds to uplift.
\end{solution}

\begin{solution}{exo:b3:quantitative-geomorphology:6}
With $x$ upstream from the outlet, $\mathrm{d}x = -\mathrm{d}d$, so $\chi(d)
= \int_d^L (A_0/c\,d'^{h})^{\theta}\,\mathrm{d}d' = A_0^{\theta} c^{-\theta}
\int_d^L d'^{-h\theta}\,\mathrm{d}d' = A_0^{\theta}c^{-\theta}(L^p -
d^p)/p$. Here $p = 0.19$, $A_0^{\theta}c^{-\theta} = 10^{2.7}/2^{0.45}$,
and $\chi = \qty{2935}{m}$.
\end{solution}

\begin{solution}{exo:b3:quantitative-geomorphology:7}
$c = 2\times10^{-6}\times(10^8)^{0.5} = \qty{0.02}{m/yr}$, \qty{20}{m/kyr};
\qty{5}{km} at that speed takes \qty{250}{kyr}. The area, and with it the
celerity, decreases upstream, so the \omterm{def:b3:quantitative-geomorphology:knickpoint}{knickpoint} slows as it climbs.
\end{solution}

\begin{solution}{exo:b3:quantitative-geomorphology:8}
$\bar e = 0.40\times600 = \qty{240}{m}$; the summits rise by $0.80\times240 =
\qty{192}{m}$; the mean elevation falls by $240 - 192 = \qty{48}{m}$.
\end{solution}

\begin{solution}{exo:b3:quantitative-geomorphology:9}
$h = h_0\ln(P_0/E)$: $0.5\ln 6 = \qty{0.90}{m}$ and $0.5\ln1.5 =
\qty{0.20}{m}$. At \qty{0.4}{mm/yr} erosion outpaces the fastest production,
$P_0$: the soil is stripped and the slope becomes bare rock, its lowering
limited by weathering.
\end{solution}

\begin{solution}{exo:b3:quantitative-geomorphology:10}
At \omterm{def:b3:quantitative-geomorphology:steady-profile}{steady state} $\mathrm{d}q/\mathrm{d}x = E$ and $q = 0$ at the divide, so
$q = E x$. Then $\kappa S = E x\,(1 - S^2/S_c^2)$, a quadratic whose positive
root is $S = [-\kappa + \sqrt{\kappa^2 + 4E^2x^2/S_c^2}\,]/(2Ex/S_c^2)$. At
20, 50, 100 and \qty{200}{m}: 0.36, 0.68, 0.89 and 1.03, against 0.4, 1.0,
2.0 and 4.0 for linear diffusion ($S = Ex/\kappa$). The hill has a rounded
crest, then sides whose gradient tends to $S_c$: nearly planar slopes, much
lower relief than the parabola.
\end{solution}

\begin{solution}{exo:b3:quantitative-geomorphology:11}
By the kinematic-wave theorem the change reaches a point of coordinate
$\chi$ after $\chi/(KA_0^m)$; the network has fully responded when it has
reached the last channel head, of largest $\chi$. $K A_0^m = 10^{-6}\times
1000 = \qty{e-3}{m/yr}$: the trunk responds in \qty{5.46}{Myr}, the network
in \qty{7.24}{Myr}, the order of a million years expected.
\end{solution}

\begin{solution}{exo:b3:quantitative-geomorphology:12}
Above the junction the tributary's steady gradient is $S = (U/2K)A^{-\theta}$,
so $\mathrm{d}z/\mathrm{d}\chi = (U/2K)A_0^{-\theta}$, half the trunk's.
The \omterm{def:b3:quantitative-geomorphology:knickpoint}{knickpoint} reaches the junction when the trunk's is at $\chi_j$; after a
further time $\Delta t$ the trunk's has moved to $\chi_k$ with $\chi_k -
\chi_j = K A_0^m\,\Delta t$, while the tributary's, with celerity $2KA^m$,
has moved to $\chi$ with $\chi - \chi_j = 2KA_0^m\,\Delta t = 2(\chi_k -
\chi_j)$.
\end{solution}

\begin{solution}{pb:b3:quantitative-geomorphology:1}
\textbf{1.} $4.5\times12000^{1.7} = \qty{3.87e7}{m^2}$, \qty{38.7}{km^2};
$4.5\times300^{1.7} = \qty{7.3e4}{m^2}$, \qty{0.073}{km^2}.

\textbf{2.} $\theta \approx -\log(0.007855/0.06641)/\log(37.08/0.388) =
0.928/1.980 = 0.47$.

\textbf{3.} Yes: with a scatter of about \qty{10}{\percent} on twelve points,
$0.466$ does not differ significantly from 0.5.

\textbf{4.} $0.06641\times\sqrt{388\,000} = \qty{41.4}{m}$;
$0.007855\times\sqrt{3.708\times10^7} = \qty{47.8}{m}$. The geometric mean,
\qty{45.1}{m}, estimates $k_s$ for $\theta = 0.5$.

\textbf{5.} A gradient is the difference of two elevations, each uncertain
by decimetres to metres, divided by a short distance; on gentle reaches the
error is a large part of the difference. $\chi$ integrates instead of
differentiating.

\textbf{6.} $p = 1 - 1.7\times0.5 = 0.15$; $\chi = (1000/\sqrt{4.5})\,
(12000^{0.15} - 9000^{0.15})/0.15 = \qty{543.1}{m}$.

\textbf{7.} $(193.4 - 170)/543.1 = 0.0431$; $k_{sn} = 0.0431\times A_0^{0.5}
= \qty{43.1}{m}$.

\textbf{8.} The chi estimate: it rests on elevations over \qty{3}{km} of
river, not on twelve noisy gradients.

\textbf{9.} $K = U/k_{sn} = 4.3\times10^{-5}/43.1 = \qty{1.00e-6}{yr^{-1}}$.

\textbf{10.} $10^{-6}\times\sqrt{3.87\times10^7}\times0.00691 =
\qty{4.30e-5}{m/yr}$, \qty{43}{m/Myr}: the outlet keeps pace with the
uplift.

\textbf{11.} At the outlet $10^{-6}\times6222 = \qty{6.2}{m/kyr}$; at the
waterfall, $d = \qty{7.40}{km}$, $A = \qty{1.70e7}{m^2}$ and
\qty{4.1}{m/kyr}.

\textbf{12.} $170 + 43\times(4.5)^{-0.5}(12000^{0.15} - 300^{0.15})/0.15 =
\qty{405}{m}$: a relief of \qty{235}{m}.

\textbf{13.} For $n = 1$ the equation is a kinematic wave: points on both
sides of the step rise by $U$ per unit time, so the step keeps its height,
the \qty{25}{m} of the fall.

\textbf{14.} $\chi = (1000/\sqrt{4.5})(12000^{0.15} - 7400^{0.15})/0.15 =
\qty{899.4}{m}$.

\textbf{15.} $t = \chi/(KA_0^{0.5}) = 899.4/10^{-3} = \qty{899}{kyr}$.

\textbf{16.} At $\chi = 899.4 + 300 = \qty{1199}{m}$:
$d^{0.15} = 12000^{0.15} - 1199\times0.15\times\sqrt{4.5}/1000$, $d =
\qty{6.25}{km}$, \qty{5.75}{km} upstream of the outlet.

\textbf{17.} At the junction ($d = \qty{9.5}{km}$) $\chi = \qty{442.8}{m}$;
the waterfall lies where the tributary's $\chi$ reaches \qty{899.4}{m}:
$s^{0.15} = 5000^{0.15} - (899.4 - 442.8)\times0.15\times\sqrt{1.5}/1000$,
$s = \qty{4.27}{km}$, \qty{0.73}{km} above the junction.

\textbf{18.} Their junctions are at $\chi = 1954$ and \qty{999}{m}, beyond
\qty{899}{m}: the \omterm{def:b3:quantitative-geomorphology:knickpoint}{knickpoint} has not yet passed them, so the trunk there is
still at its old level.

\textbf{19.} $E = 5\times10^{-3}\times0.0086 = \qty{4.3e-5}{m/yr}$,
\qty{43}{m/Myr}: the hills above the waterfall erode at the uplift rate, in
\omterm{def:b3:quantitative-geomorphology:steady-profile}{steady state}.

\textbf{20.} $150^2/(5\times10^{-3}) = \qty{4.5}{Myr}$, five times the age of
the waterfall.

\textbf{21.} Their foot was lowered \qty{25}{m} with the channel: their
lower parts are oversteepened, perhaps sliding, while their crests still
erode at the old rate. A wave of steepening climbs each slope; after a few
million years they return to the same parabola, \qty{25}{m} lower.

\textbf{22.} $t = \chi k_{sn}/(U A_0^{0.5})$ varies as $1/U$: from
$899/1.2 = 749$ to $899/0.8 = \qty{1124}{kyr}$.

\textbf{23.} With $n = 2$ the celerity grows with the gradient: the steep
face would travel faster than the reaches around it and stay sharp, but
its speed and the age read from it would differ.

\textbf{24.} The \omterm{def:b3:quantitative-geomorphology:knickpoint}{base level} of the Varden river fell
\textbf{\qty[text-series-to-math=true]{900}{kyr}} ago (\qty{899}{kyr} by the
calculation).
\end{solution}
